\documentclass[a4paper,11pt]{article}
\usepackage[utf8]{inputenc}
\usepackage{graphicx}
\usepackage{amsmath,amssymb}
\usepackage{booktabs}
\usepackage{longtable}
\usepackage{hyperref}
\usepackage{geometry}
\usepackage{array}
\geometry{margin=1in}

\title{\textbf{Conditional modal branch selection in physics-informed active control\\
of lid-driven cavity flow: identification, reproducibility,\\
and independent numerical verification}}

\author{Beniaiche Ahmed\(^{1}\), Larabi Abderrahim\(^{1}\), Belkadi Mbarek\(^{2}\) \\[4pt]
\(^{1}\)Laboratory of Fluid Mechanics, Ecole Militaire Polytechnique, Algiers, Algeria \\
\(^{2}\)Laboratory of Turbomachinery, Ecole Militaire Polytechnique, Algiers, Algeria}
\date{}

\begin{document}

\maketitle

\begin{abstract}
Active flow control of the lid-driven cavity is investigated using a parametric
physics-informed neural network (PINN) combined with a truncated spatio-temporal
Fourier actuation basis. Rather than assuming that optimization convergence
identifies a unique optimal control, we quantify the conditional selection,
reproducibility, and numerical realizability of low-dimensional control branches.
Over Reynolds numbers \(100\le Re\le 1000\) and prescribed dimensionless actuation
energies \(E^*\), independent trainings reveal four empirical classes: \(B_1\)
(stationary second spatial harmonic), \(B_2\) (time-dependent actuation),
\(B_{\mathrm{first}}\) (stationary first harmonic), and \(B_3\) (mixed). Their
frequencies depend on \(Re\), \(E^*\), initialization, loss
weights, and control representation. Changing the variance-regularization weight
redistributes branch accessibility without altering the Navier--Stokes problem;
removing it selects the stationary first-harmonic branch \(B_{\mathrm{first}}\)\,
whereas the second harmonic prevails at the reference weight. Seed experiments at
fixed \((Re,E^*)\) yield distinct converged controls with
similar losses. Selected controls are transferred without retraining to an
independent multiple-relaxation-time lattice-Boltzmann (MRT-LBM) solver, validated
against Ghia et al.; the second-harmonic control reduces integrated viscous
dissipation by \(76.5\%\) at \(Re=500\) and \(82.3\%\) at \(Re=1000\) relative to a
uniform lid, with grid-convergence analysis bounding the discretization
uncertainty. A focused iso-energetic comparison shows that this frequently
selected branch is actually more dissipative than a uniform lid at the same energy
budget (dissipation ratio in \([1.072,1.188]\)), while the first harmonic is
consistently less dissipative (ratio in \([0.646,0.726]\)), ranking first-harmonic
$<$ uniform $<$ second-harmonic at both resolutions. Reproducibility experiments
yield bit-identical double runs, and multi-seed audits confirm the branch taxonomy
(0/185 reclassification disagreements in the reference campaign at
\(\lambda_{\mathrm{var}}=1\)); the results support conditional branch
accessibility and reproducibility but refute uniqueness or global optimality.

\medskip\noindent\textbf{Keywords:} physics-informed neural networks; active flow
control; lid-driven cavity; modal branch selection; lattice-Boltzmann method;
reproducibility.
\end{abstract}

%--------------------------------------------------------------
\section*{List of Abbreviations and Variables}
%--------------------------------------------------------------
\begin{table}[h]
\centering\small
\begin{tabular}{ll}
\toprule
\textbf{Symbol} & \textbf{Description}\\
\midrule
PINN & Physics-informed neural network\\
LBM  & Lattice-Boltzmann method\\
MRT  & Multiple-relaxation-time\\
DNS  & Direct numerical simulation\\
PDE  & Partial differential equation\\
D2Q9 & Two-dimensional nine-velocity lattice\\
LOCO & Leave-one-Reynolds-condition-out\\
SR   & Symbolic regression\\
\(B_1,B_2,B_{\mathrm{first}},B_3\) & Empirical branch classes\\
\(f_1\) & Stationary first-harmonic control-energy fraction\\
\(f_2\) & Stationary second-harmonic control-energy fraction\\
\(f_t\) & Explicitly time-dependent control-energy fraction\\
\(H_N\) & Normalized branch-selection entropy\\
\(Re\) & Reynolds number\\
\(E^*\) & Prescribed target actuation energy\\
\(E_r\) & Realized actuation energy\\
\(\lambda_{\mathrm{var}}\) & Variance-regularization weight\\
\(AR\) & Cavity aspect ratio (\(L_x/L_y\))\\
\(\varepsilon\) & Integrated viscous dissipation (LBM)\\
\(\omega_{\mathrm{SR}}\) & Symbolic-regression spatial frequency\\
\(R_K\) & Ratio of fluctuating to mean kinetic energy\\
GCI  & Grid-convergence index\\
RMSE & Root-mean-square error\\
\bottomrule
\end{tabular}
\end{table}

%--------------------------------------------------------------
\section{INTRODUCTION}
\label{sec:intro}
%--------------------------------------------------------------

Active flow control seeks to modify fluid motion through externally imposed
actuation in order to achieve a prescribed dynamical or energetic objective,
including drag and lift modification, separation delay, mixing enhancement, and
viscous dissipation reduction~\cite{gad2000flow,brunton2015closed}. The difficulty
of this problem arises from the nonlinear and high-dimensional character of the
Navier--Stokes equations: even when the actuator is parameterized by a small number
of degrees of freedom, its interaction with the flow produces multiple competing
responses, strongly non-convex optimization landscapes, and, in many configurations,
genuinely multiple stationary states. Reviews of flow control and machine-learning-based
fluid mechanics have consequently emphasized the importance of reduced representations,
interpretable flow structures, and systematic treatment of nonlinear effects rather
than relying solely on scalar performance indicators
\cite{brunton2015closed,gad2000flow,duraisamy2019turbulence,brenner2019perspective}.
Methodological rigor in reporting convergence outcomes is especially important
when physics-informed or data-driven methods are used, since a single converged
solution may not represent the full landscape of accessible outcomes, and its
physical significance can only be assessed by characterizing how robustly it is
reached and whether it persists under independent numerical scrutiny.

Low-dimensional representations provide a natural route to tractable flow-control
problems. Proper orthogonal decomposition, balanced truncation, dynamic mode
decomposition, resolvent analysis, and related techniques have become standard
tools for identifying dynamically important flow structures and for constructing
reduced-order models amenable to analysis and control
\cite{rowley2017model,taira2017modal,schmid2010dmd,noack2003hierarchy}. In control
applications, such representations substantially compress the optimization problem
while retaining structures relevant to the dominant flow response. Reduced-order
optimal-control formulations have been applied to direct numerical simulation (DNS)-based flow-control calculations
\cite{bergmann2005optimal,bewley2001dns}, and reinforcement-learning methods have
extended this framework to complex geometries including bluff-body drag reduction
\cite{rabault2019drl,garnier2021drl}. The lid-driven cavity occupies a canonical
position in these investigations: its geometry and boundary conditions are simple,
yet its recirculating flow structure and Reynolds-number dependence provide a
non-trivial test of numerical and physical hypotheses, and high-quality benchmark
solutions are available for independent verification \cite{ghia1982high}. At
moderate and high Reynolds numbers the uncontrolled cavity admits multiple
topological flow states, reinforcing the need for careful interpretation of any
optimization outcome.

Physics-informed neural networks (PINNs) offer an alternative route to
reduced representations~\cite{raissi2019pinn,cai2021pinnreview}. Rather than
learning a flow response exclusively from precomputed simulation data, PINNs
incorporate the governing differential equations, boundary conditions, and physical
constraints directly into the optimization objective through automatic differentiation
of the network residuals. This formulation is particularly attractive for control
identification because the flow field and the actuation parameters are optimized
simultaneously, allowing the control itself to become an object of analysis rather
than merely an output of a separate surrogate model. Machine learning has been
applied broadly to fluid mechanics for model reduction, turbulence closure, and
optimization~\cite{duraisamy2019turbulence,brenner2019perspective,lozano2020machine},
and physics-informed variants have demonstrated promising results for laminar,
transitional, and turbulent flows. However, the physics-informed formulation does
not remove the underlying non-convexity of the optimization problem. Different
random initializations, network capacities, loss-weight configurations, regularization
terms, or control parameterizations can lead to qualitatively different converged
states. The repeated appearance of a particular low-dimensional structure cannot,
by itself, establish that the structure is unique, globally optimal, or independent
of the chosen representation.

This issue motivates a distinction that is rarely made explicit in physics-informed
flow-control studies: the distinction between the \emph{existence} of a controlled
branch, its \emph{accessibility} to a particular optimization procedure, and its
\emph{global optimality}. A control branch may be physically realizable and
repeatedly recovered while remaining only one member of a set of competing
converged states. Conversely, a branch that is infrequently recovered may not
be physically less relevant; it may simply reside in a smaller basin of attraction
under the adopted optimization formulation. The control parameterization introduces
a further layer of conditioning: a finite basis does not merely reduce the number
of degrees of freedom, but defines the subspace of control space in which the
optimizer searches. An apparently low-dimensional control structure can therefore
reflect both the dynamics of the controlled flow and the geometry of the chosen
optimization space~\cite{taira2017modal,bagheri2009global,sipp2010global}.
Subcritical non-convex landscapes in fluid mechanics have been characterized for
variational and adjoint-based methods~\cite{farano2016subcritical}, and resolvent
analysis has shown how structural assumptions can influence perceived optimal
structures~\cite{beneddine2016conditions}. A quantitative framework for
distinguishing these contributions is therefore essential before drawing physical
conclusions from a repeatedly observed modal structure.

The present study addresses this branch-selection problem for distributed active
boundary control of the two-dimensional lid-driven cavity. We formulate a
physics-informed optimization in which the Navier--Stokes equations constrain
the flow while a prescribed actuation-energy level and a dissipation-related
objective determine the control. The actuation is represented in a truncated
spatio-temporal Fourier basis, and a parametric neural representation identifies
controls over the investigated Reynolds-number interval. Rather than treating the
dominant modal structure as an assumed optimum, we explicitly test its
recoverability, multiplicity, and parameter dependence. Independent training
realizations estimate empirical branch-selection frequencies
\begin{equation}
\widehat P\!\left(B_i\mid Re,E^*,\lambda_{\mathrm{var}}\right),
\label{eq:empirical_freq}
\end{equation}
while systematic variations of \(Re\), \(E^*\), network capacity, control-space
truncation, cavity aspect ratio, loss weighting, and control basis determine when
the observed branches persist or change. The resulting branch structure is then
subjected to an independent numerical test: selected controls are transferred
without retraining to an independent multiple-relaxation-time lattice-Boltzmann
(MRT-LBM) solver, providing a completely independent discretization and
implementation. Grid refinement and benchmark comparisons quantify the associated
numerical uncertainty, and symbolic regression provides a compact analytical
approximation enabling assessment of cross-condition generalization. The objective
is not to claim the discovery of a new universal cavity mode or to establish global
optimality, but to provide a quantitative framework for distinguishing branch
existence, branch accessibility, and branch reproducibility in physics-informed
flow-control optimization.

The central questions of this study are therefore: (i) Does a distributed
physics-informed control problem repeatedly select a low-dimensional modal branch
over a range of flow and actuation parameters? (ii) Do independent initializations
reveal competing branches with comparable optimization performance? (iii) How does
the variational structure of the loss redistribute the empirical accessibility of
those branches? (iv) Does the observed modal structure survive changes in the
control representation? (v) Can the resulting branches be independently realized
and distinguished by a conventional numerical flow solver? Answering these questions
provides a more stringent interpretation of low-dimensional structures obtained from
physics-informed optimization and places modal concentration within the broader
framework of nonlinear branch selection in controlled fluid flows. The remainder of
the paper is organized as follows. Section~\ref{sec:problem_and_pinn} formulates the
control problem and PINN framework. Section~\ref{sec:global_landscape} presents the
global branch landscape. Section~\ref{sec:accessibility} investigates conditional
accessibility. Section~\ref{sec:robustness} examines robustness to capacity,
dimension, geometry, and basis. Section~\ref{sec:lbm} presents the independent LBM
verification. Section~\ref{sec:discussion} synthesizes the results.
Section~\ref{sec:sr} presents the symbolic representation.
Section~\ref{sec:conclusions} states the conclusions.

%--------------------------------------------------------------
\section{PROBLEM FORMULATION AND PHYSICS-INFORMED OPTIMIZATION}
\label{sec:problem_and_pinn}
%--------------------------------------------------------------

The objective is to examine how a distributed boundary-control problem can admit
competing low-dimensional actuation branches and how their accessibility changes
with the physical and variational parameters of the identification procedure.
Before introducing the physics-informed optimization, the controlled flow problem
and the admissible actuation space must be specified independently of the neural
network representation. All notation is consistent throughout the paper; additional
implementation details are provided in the Supplementary Material.

\subsection{Controlled cavity and governing equations}

We consider a two-dimensional incompressible lid-driven cavity in which the upper
boundary is subjected to a spatially and temporally distributed tangential actuation.
The lower and lateral walls remain stationary with no-slip conditions. The lid
velocity is not prescribed \emph{a priori} as a uniform profile but is treated as
a control variable drawn from a finite-dimensional admissible space. This is
essential for the subsequent branch analysis: the branches investigated here are
not solutions of different physical problems, but different control-flow realizations
admitted by the same governing equations, boundary geometry, and actuation constraints.

After nondimensionalization by the cavity height \(L_y^{\mathrm{dim}}\), the
wall-normal extent is unity and the streamwise extent equals the aspect ratio
\(AR=L_x^{\mathrm{dim}}/L_y^{\mathrm{dim}}\):
\begin{equation}
0\le x\le L_x,\quad 0\le y\le 1,\quad L_x=AR.
\label{eq:domain}
\end{equation}
The reference square cavity corresponds to \(L_x=1\). Velocities are
nondimensionalized by a reference lid velocity \(U_0\), time by
\(L_y^{\mathrm{dim}}/U_0\), and pressure by \(\rho U_0^2\). The Reynolds number is
\begin{equation}
Re=\frac{U_0\,L_y^{\mathrm{dim}}}{\nu},
\label{eq:reynolds}
\end{equation}
where \(\nu\) is the kinematic viscosity. Note that \(U_0\) is a scaling velocity,
not a requirement that the lid remain uniform; the distributed control
\(U_{\mathrm{lid}}\) may change sign locally provided it remains within the
admissible control space and satisfies the prescribed energy constraint.

The controlled velocity field \(\mathbf{u}=(u,v)\) and kinematic pressure \(p\)
satisfy the dimensionless incompressible Navier--Stokes equations
\begin{equation}
\frac{\partial\mathbf{u}}{\partial t}+(\mathbf{u}\cdot\nabla)\mathbf{u}
=-\nabla p+\frac{1}{Re}\nabla^2\mathbf{u},\quad
\nabla\cdot\mathbf{u}=0.
\label{eq:NS}
\end{equation}
The stationary boundaries satisfy the no-slip condition
\begin{equation}
\mathbf{u}=\mathbf{0},\quad x=0,\; x=L_x,\; y=0,
\label{eq:stationary_bc}
\end{equation}
and the controlled upper boundary satisfies
\begin{equation}
u(x,1,t;Re)=U_{\mathrm{lid}}(x,t;Re),\quad v(x,1,t;Re)=0.
\label{eq:lid_bc}
\end{equation}
The problem is nonlinear for two reasons: the Navier--Stokes convective term, and
the boundary condition~\eqref{eq:lid_bc}, which is itself selected from a
parameterized control space. The optimization therefore couples the selection of a
boundary actuation to the nonlinear flow response that the same actuation generates.

\subsection{Finite-dimensional spatio-temporal control}

The lid actuation is represented in a truncated spatio-temporal Fourier basis.
The sine basis in \(x\) enforces \(U_{\mathrm{lid}}(0,t;Re)=U_{\mathrm{lid}}(L_x,t;Re)=0\),
preventing discontinuous tangential velocity at the upper corners for every
tested aspect ratio. The cosine basis in \(t\) allows the \(n=0\) terms to
represent purely stationary actuation, the most directly interpretable patterns
physically. The lid velocity is expanded as
\begin{equation}
U_{\mathrm{lid}}(x,t,Re)=\sum_{m=1}^{N_x}\sum_{n=0}^{N_t-1}
a_{mn}(Re)\sin\!\left(\frac{m\pi x}{L_x}\right)\cos(n\pi t).
\label{eq:fourier_control}
\end{equation}
All modal labels use the physical (one-based) convention: the dominant stationary
second spatial harmonic corresponds to \((m,n)=(2,0)\). The baseline control space
uses \((N_x,N_t)=(6,5)\), giving 30 nominal coefficients. Control-space sensitivity
examines
\begin{equation}
(N_x,N_t)\in\{(4,3),(6,1),(6,3),(6,5),(8,5),(8,9)\},
\label{eq:truncation_sweep}
\end{equation}
spanning 12 to 72 nominal modes. The modal energy of coefficient \((m,n)\) is
\begin{equation}
\mathcal{E}_{m0}=\tfrac{1}{2}a_{m0}^2,\quad
\mathcal{E}_{mn}=\tfrac{1}{4}a_{mn}^2\quad(n\ge1),
\label{eq:modal_energy}
\end{equation}
where the factors account for temporal basis normalization
(\(\langle\cos^2(0)\rangle=1\), \(\langle\cos^2(n\pi t)\rangle=\tfrac{1}{2}\) for \(n\ge1\)).
The normalized modal fraction is
\begin{equation}
f_{mn}=\frac{\mathcal{E}_{mn}}{\displaystyle\sum_{i=1}^{N_x}\sum_{j=0}^{N_t-1}\mathcal{E}_{ij}},
\qquad\sum_{m=1}^{N_x}\sum_{n=0}^{N_t-1}f_{mn}=1.
\label{eq:modal_fraction}
\end{equation}
The stationary second-harmonic fraction is \(f_2\equiv f_{20}\), and the total
time-dependent fraction is \(f_t=\sum_{m=1}^{N_x}\sum_{n=1}^{N_t-1}f_{mn}\).
The realized actuation energy is
\begin{equation}
E_r=\langle U_{\mathrm{lid}}^2\rangle_{s,t}
=\frac{1}{T}\int_0^T\!\!\int_0^1 U_{\mathrm{lid}}^2(s,t;Re)\,ds\,dt,
\quad s=\frac{x}{L_x},
\label{eq:realized_energy}
\end{equation}
and the optimization is conditioned on \(E_r\approx E^*\). Modal-energy quadrature
conventions are detailed in Supplementary Sec.~S2.

\subsection{Parametric PINN formulation}

The flow field and Reynolds-dependent control coefficients are represented
simultaneously by a parametric neural network with inputs \((x,y,t,\widetilde{Re})\),
where
\begin{equation}
\widetilde{Re}=\frac{Re-550}{450}\in[-1,1],
\label{eq:re_norm}
\end{equation}
maps the training interval \([100,1000]\) to \([-1,1]\). A single trained network
represents the flow-control pair over the entire Reynolds-number range, enabling
efficient evaluation of the branch map without repeated retraining at each condition.
The flow variables \((u,v,p)\) and control coefficients \(\{a_{mn}(\widetilde{Re})\}\)
are the network outputs. Two subnetworks are employed: a flow network with 6 hidden
layers of 64 neurons each using \(\tanh\) activation (21\,379 trainable parameters),
and a control network with 4 hidden layers of 32 neurons each (3\,614 parameters).
All computations are performed in double-precision arithmetic, and all physical
residuals are evaluated by automatic differentiation. Full architecture
specifications are provided in Supplementary Sec.~S1.

\subsection{Optimization objective and training procedure}

The multi-objective loss combines partial-differential-equation (PDE),
boundary-condition, dissipation, control-energy, variance, and
Reynolds-separation terms:
\begin{equation}
\mathcal{L}=\lambda_{\mathrm{PDE}}\mathcal{L}_{\mathrm{PDE}}
+\lambda_{\mathrm{BC}}\mathcal{L}_{\mathrm{BC}}
+\lambda_{\mathrm{diss}}\mathcal{L}_{\mathrm{diss}}
+\lambda_{\mathrm{ctrl}}\mathcal{L}_{\mathrm{ctrl}}
+\lambda_{\mathrm{var}}\mathcal{L}_{\mathrm{var}}
+\lambda_{Re}\mathcal{L}_{Re},
\label{eq:total_loss}
\end{equation}
with reference weights
\(\lambda_{\mathrm{PDE}}=1,\; \lambda_{\mathrm{BC}}=15,\;
\lambda_{\mathrm{diss}}=0.05,\; \lambda_{\mathrm{ctrl}}=200,\;
\lambda_{\mathrm{var}}=1,\; \lambda_{Re}=2\).
The individual loss components are:

\textbf{PDE residual loss:}
\(\mathcal{L}_{\mathrm{PDE}}=\langle R_u^2+R_v^2+R_c^2\rangle_{\mathcal{X}_{\mathrm{int}}}\),
where \(R_u,R_v,R_c\) are momentum and continuity residuals evaluated at
\(N_{\mathrm{phys}}=4\,000\) interior collocation points resampled uniformly at each epoch.

\textbf{Boundary condition loss:}
\(\mathcal{L}_{\mathrm{BC}}=\langle u^2+v^2\rangle_{\Gamma_{\mathrm{stat}}}
+\langle[u-U_{\mathrm{lid}}]^2+v^2\rangle_{\Gamma_{\mathrm{top}}}\),
where \(\Gamma_{\mathrm{stat}}=\{x=0\}\cup\{x=L_x\}\cup\{y=0\}\). The weight
\(\lambda_{\mathrm{BC}}=15\) enforces boundary conditions strongly relative to
the interior residual.

\textbf{Viscous dissipation objective:}
\begin{equation}
\mathcal{L}_{\mathrm{diss}}=\int_\Omega\frac{2}{Re}
\!\left[u_x^2+v_y^2+\tfrac{1}{2}(u_y+v_x)^2\right]d\Omega.
\label{eq:dissipation}
\end{equation}
This term drives the optimizer toward low-dissipation actuation patterns and
constitutes the physical objective of the control problem.

\textbf{Actuation energy penalty:}
\(\mathcal{L}_{\mathrm{ctrl}}=(E_r-E^*)^2\). The large weight
\(\lambda_{\mathrm{ctrl}}=200\) enforces the energy budget tightly, ensuring that
competing converged solutions are compared at the same actuation energy level and
that the modal competition is not trivially explained by different energy budgets.

\textbf{Variance regularization:}
\begin{equation}
\mathcal{L}_{\mathrm{var}}=-\log\!\left[\operatorname{Var}(U_{\mathrm{lid}})+10^{-5}\right].
\label{eq:var_loss}
\end{equation}
This term penalizes spatially uniform or near-uniform lid profiles by rewarding
higher spatial variance. It plays a critical role in breaking degeneracy toward
trivially flat controls and is the primary determinant of branch selection, as
demonstrated in Section~\ref{sec:accessibility}.

\textbf{Reynolds-separation regularization:}
\begin{equation}
\mathcal{L}_{Re}=\max\!\left[0,\;\delta-
\frac{1}{N_{\mathrm{probe}}}\sum_{i=1}^{N_{\mathrm{probe}}}
\!\left(U_{\mathrm{lid}}(x_i,t_0,Re_{\min})
-U_{\mathrm{lid}}(x_i,t_0,Re_{\max})\right)^2\right],
\end{equation}
with \(N_{\mathrm{probe}}=20\), \(t_0=0.5\), \(\delta=0.01\), preventing the
parametric network from collapsing to a Reynolds-independent control.

Training consists of a 300-epoch pre-training phase in which only the flow network
is updated with control coefficients frozen at zero, establishing an initial
physically consistent flow representation. A 3\,000-epoch joint optimization phase
follows using the Adam optimizer with initial learning rate \(\eta_0=8\times10^{-4}\)
and exponential decay by factor \(0.995\) per 100 epochs. Collocation points are
resampled at every epoch. Complete sampling details are provided in
Supplementary Sec.~S1.

\subsection{Branch classification and reproducibility protocol}

Converged neural realizations are classified into four empirical branch classes
according to their modal signature computed from
Eqs.~\eqref{eq:modal_energy}--\eqref{eq:modal_fraction}. Let
\(f_1\equiv f_{10}\) denote the stationary first-harmonic (fundamental) fraction,
\(f_2\equiv f_{20}\) the stationary second-harmonic fraction, and
\(f_t=\sum_m\sum_{n\ge1}f_{mn}\) the total explicitly time-dependent fraction:
\begin{itemize}
\item \textbf{Branch \(B_{\mathrm{first}}\) --- Stationary first-harmonic dominated:}
\(f_1>0.5\). Physical harmonic \((m,n)=(1,0)\).
\item \textbf{Branch \(B_1\) --- Stationary second-mode dominated:}
\(f_2>2f_t\) and \(f_2>0.5\).
Physical harmonic \((m,n)=(2,0)\). Ensemble mean \(\overline{f_2}=85.3\%\),
\(\overline{f_t}=7.4\%\).
\item \textbf{Branch \(B_2\) --- Temporally dominated:} \(f_t>2f_2\) and \(f_t>0.5\).
Physical harmonic \((m,n)=(1,1)\). Ensemble mean \(\overline{f_2}=7.5\%\),
\(\overline{f_t}=88.1\%\).
\item \textbf{Branch \(B_3\) --- Mixed:} none of the above.
Ensemble mean \(\overline{f_2}=34.8\%\), \(\overline{f_t}=43.9\%\).
\end{itemize}
The empirical selection frequency over \(N_s\) independent seeds is
\(\widehat P(B_k)=N_k/N_s\). Selection diversity is measured by normalized Shannon
entropy \(H_N=-\frac{1}{\ln 4}\sum_k\widehat P(B_k)\ln\widehat P(B_k)\), which
equals zero for deterministic single-branch selection and unity for a uniform
distribution over all four classes. The reference campaign of
Sec.~\ref{sec:global_landscape}, conducted at \(\lambda_{\mathrm{var}}=1\), contains
no \(B_{\mathrm{first}}\) realization; this branch becomes populated only when the
variance regularizer is weakened or removed (Sec.~\ref{sec:lambda_var}). In the
original three-class scheme these first-harmonic realizations were grouped with the
mixed branch \(B_3\); the explicit four-class classification shows that they
constitute a distinct attractor. Detailed implementation parameters and
realization-level campaign data are provided in the Supplementary Material.

%--------------------------------------------------------------
\section{GLOBAL BRANCH-SELECTION LANDSCAPE}
\label{sec:global_landscape}
%--------------------------------------------------------------

\subsection{Broad \((Re,E^*)\) campaign}

The broad campaign evaluated 185 independent training runs across 30 parameter-space
cells (\(Re\in\{100,300,500,700,900,1000\}\), \(E^*\in\{0.1,0.25,0.5,1,2\}\)),
with at least 5 seeds per cell and up to 20 at selected cells of particular
interest. Figure~\ref{fig:branch_selection_map} presents the resulting empirical
branch-frequency map \(\widehat P(B_k)\), and Table~\ref{tab:branch_statistics}
summarizes the global realization statistics. The complete realization-level dataset
is provided in Supplementary Sec.~S3 and Supplementary Table~S1.

\begin{figure*}[t]
\centering
\includegraphics[width=0.92\textwidth]{fig01_branch_selection_map.png}
\caption{Empirical branch-selection frequencies over the \((Re,E^*)\) parameter
space from the 185-realization campaign. Each panel reports the fraction of converged
realizations classified as \(B_1\) (stationary second-harmonic dominated), \(B_2\)
(temporally dominated), or \(B_3\) (mixed) at each parameter cell. The values
represent empirical recovery frequencies from independent optimizations and should not
be interpreted as deterministic phase boundaries. Cell-level realization counts and
realization-level classifications are reported in Supplementary Table~S1.}
\label{fig:branch_selection_map}
\end{figure*}

\begin{table}[t]
\centering
\caption{Global statistics of the empirical branch classes populated over the
185-realization \((Re,E^*)\) campaign at \(\lambda_{\mathrm{var}}=1\) (no
\(B_{\mathrm{first}}\) realization occurs at this weight). Modal fractions in
percentages.}
\label{tab:branch_statistics}
\small
\begin{tabular}{lcccccc}
\toprule
Branch & \(N\) & Frequency (\%) & \(\overline{f_2}\) (\%) & Median \(f_2\) (\%) & \(\overline{f_t}\) (\%) & Median \(f_t\) (\%)\\
\midrule
\(B_1\) & 78 & 42.2 & 85.3 & 88.9 & 7.4  & 3.8\\
\(B_2\) & 81 & 43.8 & 7.5  & 3.5  & 88.1 & 95.4\\
\(B_3\) & 26 & 14.1 & 34.8 & 36.3 & 43.9 & 45.1\\
\bottomrule
\end{tabular}
\end{table}

Across the 185 runs, \(B_1\), \(B_2\), and \(B_3\) account for 42.2\%, 43.8\%,
and 14.1\% of solutions, respectively. The near balance between \(B_1\) and \(B_2\)
is the most significant feature of the global landscape: while the stationary
second-harmonic branch is frequently accessible, it does not dominate the ensemble
to the point of being the unique or near-deterministic outcome. The mixed branch
\(B_3\) is the least frequently selected, arising primarily at elevated actuation
energies \(E^*\ge1\) and higher Reynolds numbers. No \(B_{\mathrm{first}}\)
realization occurs anywhere in this campaign: the first-harmonic fraction remains
below \(f_1=0.13\) throughout, confirming that the reference weight
\(\lambda_{\mathrm{var}}=1\) actively excludes the fundamental mode. The branch
frequencies exhibit
a clear \(Re\)-dependence: \(B_1\) is more readily accessible at moderate \(Re\)
around 500, while at \(Re\in[100,300]\), \(B_2\) often dominates. This suggests
that the relative depths of the two competing optimization basins vary systematically
with the flow regime.

\subsection{Continuous modal structure in observable space}

Figure~\ref{fig:branch_observable_space} illustrates the distribution of all 185
realizations in the two-dimensional \((f_2,f_t)\) observable space. Three distinct
clusters are clearly visible, separated by low-density gaps. The \(B_1\) cluster
concentrates near the \(f_t\approx0\) axis at high \(f_2\); the \(B_2\) cluster
concentrates near the \(f_2\approx0\) axis at high \(f_t\); and the \(B_3\) cluster
occupies an intermediate region. The cluster separation establishes that the
branch classes are well-defined structural entities rather than arbitrary divisions
of a continuous distribution. The first-harmonic class \(B_{\mathrm{first}}\) does
not appear in this observable space because it is not populated at
\(\lambda_{\mathrm{var}}=1\); it emerges only for weakened variance regularization
(Sec.~\ref{sec:lambda_var}).

\begin{figure}[t]
\centering
\includegraphics[width=0.82\linewidth]{fig02_branch_observable_space.png}
\caption{Distribution of the 185 converged realizations in the \((f_2,f_t)\) modal
observable space, colored by assigned empirical branch class (\(B_1\): red;
\(B_2\): blue; \(B_3\): green). Three well-separated clusters confirm that the
empirical branch classes correspond to structurally distinct attractors of the
physics-informed optimization.}
\label{fig:branch_observable_space}
\end{figure}

\subsection{Conditional branch accessibility and non-determinism}

A key observation from Figure~\ref{fig:branch_selection_map} is that \((Re,E^*)\)
does not uniquely determine a single control class. At \(Re=500\), \(E^*=0.25\),
a 15-seed sub-campaign yields \(\widehat P(B_1)=0.467\), \(\widehat P(B_2)=0.400\),
\(\widehat P(B_3)=0.133\), demonstrating genuine multi-branch accessibility at a
single operating point. At \(Re=1000\), \(E^*=1\), a 20-seed sub-campaign yields
\(\widehat P(B_1)=0.55\), \(\widehat P(B_2)=0.20\), \(\widehat P(B_3)=0.25\).
The mapping \((Re,E^*)\to\mathcal{B}\) is therefore set-valued: each operating
point maps to an empirical distribution \(\{\widehat P(B_1),\widehat P(B_2),
\widehat P(B_3)\}\). The following sections investigate the conditions under
which this distribution changes.

Finally, the temporal structure of converged realizations was examined at the
trajectory level for seeds representative of the three classes (seed 7
\(\to B_1\), seed 8 \(\to B_2\), seed 4 \(\to B_3\)). A clear four-epoch
spectral peak in the raw per-epoch training loss (period \(\approx4\) epochs;
power up to \(\sim55\times\) the spectral median) does \emph{not} propagate to
the control energy, the control coefficients, or the modal fractions: no
period-4 signature is found in the control path at the sampling scale, and the
branch labels remain stable across transients. The four-epoch motif is therefore
an optimizer-internal oscillation confined to the loss trajectory, not a physical
or control-space cycle. The three classes differ in their established behavior:
the \(B_1\) control is effectively frozen (control displacement
\(D_1\simeq0.009\)), the \(B_2\) realization shows a finite relocation transient
after which the control freezes (stationary \(D_1\simeq0.008\)), and the \(B_3\)
realization exhibits a slow, persistent drift of both the control coefficients and
the modal fractions (stationary modal displacement \(\Delta f(4)\simeq0.05\),
about \(50\times\) the \(B_1\) value). Autocorrelation functions, cycle tests,
and their definitions are provided in Supplementary Sec.~S11.

%--------------------------------------------------------------
\section{CONDITIONAL BRANCH ACCESSIBILITY}
\label{sec:accessibility}
%--------------------------------------------------------------

\subsection{Variational dependence on \(\lambda_{\mathrm{var}}\)}
\label{sec:lambda_var}

The effect of variance regularization was tested in a dedicated 90-run campaign
(\(N_s=5\) seeds at 3 operating points, 6 values of
\(\lambda_{\mathrm{var}}\in\{0,0.25,0.5,1,2,5\}\)).
Table~\ref{tab:lambda_branch_selection} and
Figure~\ref{fig:lambda_branch_probabilities} demonstrate that changing
\(\lambda_{\mathrm{var}}\) systematically redistributes empirical branch-selection
frequencies while leaving the Navier--Stokes problem entirely unchanged. Complete
realization-level results and finite-sample Wilson-score confidence intervals are
given in Supplementary Sec.~S4.

\begin{table*}[t]
\centering
\caption{Empirical branch-selection frequencies from the \(\lambda_{\mathrm{var}}\)
sensitivity campaign. Each entry is based on five independent realizations at
fixed \((Re,E^*)\). \(H_N\) is the normalized branch-selection entropy
(normalised by \(\ln 4\)). Complete
95\% confidence intervals are provided in Supplementary Sec.~S4.}
\label{tab:lambda_branch_selection}
\scriptsize
\begin{tabular}{cccccccccc}
\toprule
\(Re\) & \(E^*\) & \(\lambda_{\mathrm{var}}\) & \(\widehat P(B_1)\) & \(\widehat P(B_2)\) & \(\widehat P(B_{\mathrm{first}})\) & \(\widehat P(B_3)\) & \(H_N\) & Dominant class\\
\midrule
500  & 0.25 & 0    & 0.0 & 0.0 & 1.0 & 0.0 & 0.000 & \(B_{\mathrm{first}}\)\\
500  & 0.25 & 0.25 & 0.6 & 0.4 & 0.0 & 0.0 & 0.485 & \(B_1\)\\
500  & 0.25 & 0.5  & 0.4 & 0.2 & 0.0 & 0.4 & 0.761 & \(B_1/B_3\)\\
500  & 0.25 & 1    & 0.4 & 0.2 & 0.0 & 0.4 & 0.761 & \(B_1/B_3\)\\
500  & 0.25 & 2    & 0.4 & 0.4 & 0.0 & 0.2 & 0.761 & \(B_1/B_2\)\\
500  & 0.25 & 5    & 0.4 & 0.6 & 0.0 & 0.0 & 0.485 & \(B_2\)\\
\midrule
700  & 2    & 0    & 0.0 & 0.0 & 1.0 & 0.0 & 0.000 & \(B_{\mathrm{first}}\)\\
700  & 2    & 0.25 & 0.0 & 0.0 & 1.0 & 0.0 & 0.000 & \(B_{\mathrm{first}}\)\\
700  & 2    & 0.5  & 0.0 & 0.6 & 0.0 & 0.4 & 0.485 & \(B_2\)\\
700  & 2    & 1    & 0.2 & 0.4 & 0.0 & 0.4 & 0.761 & \(B_2/B_3\)\\
700  & 2    & 2    & 0.6 & 0.2 & 0.0 & 0.2 & 0.685 & \(B_1\)\\
700  & 2    & 5    & 0.6 & 0.2 & 0.0 & 0.2 & 0.685 & \(B_1\)\\
\midrule
1000 & 1    & 0    & 0.0 & 0.0 & 1.0 & 0.0 & 0.000 & \(B_{\mathrm{first}}\)\\
1000 & 1    & 0.25 & 0.0 & 0.6 & 0.0 & 0.4 & 0.485 & \(B_2\)\\
1000 & 1    & 0.5  & 0.4 & 0.4 & 0.0 & 0.2 & 0.761 & \(B_1/B_2\)\\
1000 & 1    & 1    & 0.4 & 0.4 & 0.0 & 0.2 & 0.761 & \(B_1/B_2\)\\
1000 & 1    & 2    & 0.4 & 0.4 & 0.0 & 0.2 & 0.761 & \(B_1/B_2\)\\
1000 & 1    & 5    & 0.4 & 0.6 & 0.0 & 0.0 & 0.485 & \(B_2\)\\
\bottomrule
\end{tabular}
\end{table*}

\begin{figure*}[t]
\centering
\includegraphics[width=0.90\textwidth]{fig03_lambda_branch_probabilities.pdf}
\caption{Empirical redistribution of branch-selection frequencies with the
variance-regularization weight \(\lambda_{\mathrm{var}}\) for three fixed operating
conditions. Heatmaps show the four classes \(B_1\), \(B_2\),
\(B_{\mathrm{first}}\), and \(B_3\); cell values are empirical frequencies
estimated from five independent realizations per \(\lambda_{\mathrm{var}}\) value.
95\% Wilson-score confidence intervals are tabulated in Supplementary Sec.~S4.}
\label{fig:lambda_branch_probabilities}
\end{figure*}

At \(\lambda_{\mathrm{var}}=0\), all three operating conditions yield exclusively
\(B_{\mathrm{first}}\) realizations (\(H_N=0\)): in the absence of spatial-variance
regularization the optimizer collapses onto the lowest-order stationary mode, the
single-lobe fundamental \(\sin(\pi x)\), which is spatially structured rather than
mixed. The suppression of \(B_{\mathrm{first}}\) by a positive
\(\lambda_{\mathrm{var}}\) is operating-point dependent: at \((Re,E^*)=(500,0.25)\)
and \((1000,1)\) the branch is absent already at the smallest tested positive weight
\(\lambda_{\mathrm{var}}=0.25\), whereas at \((700,2)\) \(B_{\mathrm{first}}\) persists
at \(\lambda_{\mathrm{var}}=0.25\) (5/5 realizations) and disappears only at
\(\lambda_{\mathrm{var}}=0.5\). The transition threshold is therefore operating-point
dependent rather than universally determined by the sign of \(\lambda_{\mathrm{var}}\).
For \(\lambda_{\mathrm{var}}\ge0.5\) the
distribution redistributes non-monotonically: at \((Re,E^*)=(500,0.25)\), \(B_1\)
frequency is maximized at intermediate \(\lambda_{\mathrm{var}}\approx0.25\), whereas
larger values shift the distribution toward \(B_2\), reaching
\(\widehat P(B_2)=0.6\) at \(\lambda_{\mathrm{var}}=5\). At \((Re,E^*)=(700,2)\),
the opposite trajectory is observed: values in the range \([2,5]\) favor \(B_1\)
even though small values favor \(B_2\) or \(B_3\). These results demonstrate that
variance regularization acts as an optimization landscape bias: it modifies which
basin of attraction the optimizer enters during training without modifying the
underlying fluid-mechanics problem. This constitutes direct evidence that the branch
structure observed in physics-informed control studies is at least partially a
property of the variational formulation rather than solely of the physical problem.

\subsection{Dependence on actuation energy}

Target-energy sweeps at fixed \(Re=500\) over
\(E^*\in\{0.01,0.05,0.1,0.25,0.5,1,2\}\) reveal that the second-harmonic purity
\(f_2\) decreases monotonically as the prescribed actuation amplitude grows,
while \(f_t\) correspondingly increases (Figure~\ref{fig:energy_sweep}). At the
lowest energies (\(E^*\le0.05\)), the control budget is insufficient to explore
temporal degrees of freedom, and the optimizer concentrates almost entirely on the
second spatial harmonic (\(f_2\approx93\)--\(94\%\)). At the reference condition
\(E^*=0.25\), \(f_2=92.2\%\) with \(f_t=1.3\%\). At \(E^*=2\), however,
\(f_2\) drops to \(74.0\%\) while \(f_t\) rises to \(21.8\%\), indicating that
at large actuation amplitudes the optimizer increasingly distributes control energy
away from the second harmonic toward additional spatial and temporal components.
This transition reflects the fact that at high energy the energy budget can be
satisfied by many different modal mixtures, making the second-harmonic basin less
uniquely attractive. Throughout the sweep the first-harmonic fraction remains small
(\(f_1\le8\%\)), so the \(B_{\mathrm{first}}\) branch is never selected at the
reference weight. Extended energy sweep datasets are provided in Supplementary Sec.~S5.

\begin{figure*}[t]
\centering
\includegraphics[width=0.88\textwidth]{fig05_energy_sweep.pdf}
\caption{Dependence of modal structure on prescribed actuation energy \(E^*\) at
\(Re=500\). Left: second-harmonic fraction \(f_2\) and time-dependent fraction
\(f_t\) across the energy sweep. Right: second-harmonic amplitude \(A_2\) and
realized energy \(E_r\) versus \(E^*\). Second-harmonic dominance is robust at
low to moderate energies but erodes at \(E^*\ge1\).}
\label{fig:energy_sweep}
\end{figure*}

\subsection{Initialization sensitivity}
\label{sec:seeds}

Ten independent random seeds at \(Re=500\), \(E^*=0.25\) produced 4 second-mode
dominated (\(B_1\)), 4 temporally dominated (\(B_2\)), and 2 intermediate (\(B_3\))
realizations. Figure~\ref{fig:seed_multiplicity} illustrates the seed-to-seed
outcomes in the \((f_2,f_t)\) space. No seed selects the first-harmonic branch
(\(f_1<0.02\) for all ten seeds), confirming that \(B_{\mathrm{first}}\) is not a
low-energy accidental outcome at the reference weight. Despite qualitatively distinct
modal structures,
total training losses are tightly clustered within \([1.371,1.394]\), a relative
spread of only 1.7\%. The \(B_1\) and \(B_2\) realizations achieve essentially
equivalent loss values, confirming that the two branches are separated by distinct
basins of comparable depth rather than by a large energy gap. This multiplicity is
the direct signature of a non-convex optimization problem with competing local
minima, and it demonstrates that the selection of a particular branch by any single
training run cannot serve as evidence of global optimality. The complete seed-level
dataset is provided in Supplementary Sec.~S5.

Machine-level reproducibility was verified independently for each of the three
branch classes through idempotent double runs at the reference condition
\(Re=500\), \(E^*=0.25\) using seeds representative of \(B_1\), \(B_2\), and
\(B_3\) (seeds 7, 8, and 4, respectively). The two runs of each pair are
bit-identical to the last decimal for seeds 7 and 8 (zero difference in \(f_2\),
\(f_t\), and total loss) and quasi-identical for seed 4 (branch class \(B_3\)
unchanged; \(|\Delta f_2|=0.012\)), and each pair coincides, after re-analysis
with the final code, with the original study realization of the same seed. The
reproducibility of the branch taxonomy is therefore established per branch class
rather than only through aggregate frequencies. Detailed tables are provided in
Supplementary Sec.~S10.

\begin{figure*}[t]
\centering
\includegraphics[width=0.88\textwidth]{fig06_seed_multiplicity.pdf}
\caption{Seed-to-seed multiplicity at \(Re=500\), \(E^*=0.25\) across 10
independent initializations. Left: modal outcome in the \((f_2,f_t)\) plane,
labeled by seed index and colored by total training loss. Right: realized
control energy \(E_r\) versus total loss. Total losses lie within \([1.371,1.394]\),
confirming that \(B_1\) and \(B_2\) correspond to competing basins of comparable
depth.}
\label{fig:seed_multiplicity}
\end{figure*}

\subsection{Loss-term ablation}
\label{sec:ablation}

Ablation experiments at \(Re=500\), \(E^*=0.25\) isolate the role of individual
loss components. Table~\ref{tab:loss_ablation} reports results for single and dual
term removals. Removing \(\mathcal{L}_{\mathrm{var}}\) alone causes \(f_2\) to
collapse from \(92.2\%\) to \(0.2\%\) while the first-harmonic fraction \(f_1\)
jumps from \(0.0\%\) to \(90.4\%\): the modal class shifts completely from the
second-harmonic branch \(B_1\) to the first-harmonic branch
\(B_{\mathrm{first}}\). Removing \(\mathcal{L}_{Re}\) or
\(\mathcal{L}_{\mathrm{diss}}\) individually leaves \(f_2\) essentially unchanged
at \(92.2\)--\(92.3\%\). Dual ablations confirm that the shift induced by removing
\(\mathcal{L}_{\mathrm{var}}\) is not amplified or modulated by simultaneously
removing other terms; in every case the outcome is \(B_{\mathrm{first}}\). Reducing
\(\lambda_{\mathrm{var}}\) by half decreases second-mode concentration modestly to
\(89.3\%\), while doubling it increases it to \(92.8\%\). These results establish
that \(\mathcal{L}_{\mathrm{var}}\) is the critical loss term for \(B_1\) selection
and, equivalently, for excluding \(B_{\mathrm{first}}\), and that its action is not
redundant with the dissipation or \(Re\)-separation terms. Additional ablation
datasets are given in Supplementary Sec.~S5.

\begin{table*}[t]
\centering
\caption{Loss-term ablation results at \(Re=500\), \(E^*=0.25\). The
first-harmonic fraction \(f_1\) isolates the modal-class shift induced by removing
\(\mathcal{L}_{\mathrm{var}}\).}
\label{tab:loss_ablation}
\small
\begin{tabular}{lccccc}
\toprule
Configuration & \(E_r\) & \(f_1\) (\%) & \(f_2\) (\%) & \(f_t\) (\%) & Class\\
\midrule
Baseline & 0.2534 & 0.0 & 92.2 & 1.3 & \(B_1\)\\
No \(\mathcal{L}_{\mathrm{var}}\) & 0.2582 & 90.4 & 0.2 & 5.1 & \(B_{\mathrm{first}}\)\\
No \(\mathcal{L}_{Re}\) & 0.2534 & 0.0 & 92.2 & 1.3 & \(B_1\)\\
No \(\mathcal{L}_{\mathrm{diss}}\) & 0.2533 & 0.0 & 92.3 & 1.2 & \(B_1\)\\
No \(\mathcal{L}_{\mathrm{var}},\mathcal{L}_{Re}\) & 0.2582 & 90.4 & 0.2 & 5.1 & \(B_{\mathrm{first}}\)\\
No \(\mathcal{L}_{\mathrm{var}},\mathcal{L}_{\mathrm{diss}}\) & 0.2578 & 90.5 & 0.3 & 5.0 & \(B_{\mathrm{first}}\)\\
\(\lambda_{\mathrm{var}}/2\) & 0.2546 & 0.2 & 89.3 & 1.5 & \(B_1\)\\
\(2\lambda_{\mathrm{var}}\) & 0.2598 & 0.1 & 92.8 & 1.0 & \(B_1\)\\
\bottomrule
\end{tabular}
\end{table*}

%--------------------------------------------------------------
\section{ROBUSTNESS AND REPRESENTATION DEPENDENCE}
\label{sec:robustness}
%--------------------------------------------------------------

\subsection{Neural-network capacity}

To assess whether the second-harmonic concentration is an artifact of the specific
network architecture, three capacity variants were evaluated at \(Re=500\),
\(E^*=0.25\): a small network (\(4\times32\)), the reference network
(\(6\times64\)), and a large network (\(8\times128\)). The second-harmonic fraction
remains in the range \(f_2\in[89.5\%,92.2\%]\) across all three configurations,
a spread of only \(2.7\%\). The fact that a network roughly 40 times smaller than
the large variant produces essentially the same modal concentration argues against
the hypothesis that the second-harmonic structure is imposed by the inductive bias
of the neural network. Detailed architecture-level results are reported in
Supplementary Sec.~S6.

\subsection{Control-space dimension}

Truncation sweeps across six control dimensions (Table~\ref{tab:truncation_sensitivity})
show persistent second-mode dominance with \(f_2\in[87.3\%,92.3\%]\).
Figure~\ref{fig:capacity_truncation} confirms that expanding the basis to 72 modes
does not eliminate the \(B_1\) branch or substantially redistribute modal energy
toward higher harmonics; additional modes remain largely unactivated. The
second-harmonic amplitude \(A_2\) varies only from \(0.668\) to \(0.684\) across
all truncations. These results support the hypothesis that the second-harmonic
structure reflects a genuine low-dissipation response of the cavity flow rather
than a numerical artifact of insufficient degrees of freedom. Truncation datasets
are provided in Supplementary Sec.~S6.

\begin{table}[t]
\centering
\caption{Control-truncation sensitivity at \(Re=500\), \(E^*=0.25\).}
\label{tab:truncation_sensitivity}
\small
\begin{tabular}{cccccc}
\toprule
\((N_x,N_t)\) & Modes & \(A_2\) & \(E_r\) & \(f_2\) (\%) & \(f_t\) (\%)\\
\midrule
(4,3) & 12 & 0.671 & 0.2541 & 87.6 & 2.1\\
(6,1) &  6 & 0.678 & 0.2533 & 89.2 & 0.0\\
(6,3) & 18 & 0.678 & 0.2778 & 89.3 & 1.8\\
(6,5) & 30 & 0.684 & 0.2533 & 92.3 & 1.3\\
(8,5) & 40 & 0.669 & 0.2475 & 87.8 & 2.0\\
(8,9) & 72 & 0.668 & 0.2775 & 87.3 & 2.4\\
\bottomrule
\end{tabular}
\end{table}

\begin{figure*}[t]
\centering
\includegraphics[width=0.88\textwidth]{fig08_capacity_truncation.pdf}
\caption{Control-truncation sensitivity at \(Re=500\), \(E^*=0.25\). Left:
\(f_2\) and \(f_t\) across six tested truncations ordered by total nominal modes.
Right: second-harmonic amplitude \(A_2\) and realized energy \(E_r\). The \(B_1\)
branch persists with \(f_2\ge87.3\%\) across all tested dimensions.}
\label{fig:capacity_truncation}
\end{figure*}

\subsection{Cavity aspect ratio}

Changing the cavity geometry (\(AR\in\{0.5,1,2\}\)) at \(Re=500\), \(E^*=0.25\)
shows that the stationary second-harmonic signature remains dominant at the
reference aspect ratio (\(f_2=92.2\%\)) and at the elongated geometry
(\(AR=2\), \(f_2=91.2\%\)), but its concentration decreases notably at \(AR=0.5\),
where \(f_2\) drops to \(69.9\%\) while \(f_t\) rises to \(26.7\%\)
(Table~\ref{tab:aspect_ratio}). The compressed geometry changes the relationship
between the harmonic structure of the Fourier basis and the natural recirculation
scales of the cavity: at \(AR=0.5\), the second harmonic has a half-wavelength of
\(0.25\), shorter than the natural recirculation scale, and may be less efficient
at redistributing the primary vortex structure. Detailed basis-construction data
are provided in Supplementary Sec.~S7.

\begin{table}[t]
\centering
\caption{Aspect-ratio dependence of the modal structure at \(Re=500\), \(E^*=0.25\).}
\label{tab:aspect_ratio}
\small
\begin{tabular}{ccccc}
\toprule
\(AR\) & \(A_2\) & \(E_r\) & \(f_2\) (\%) & \(f_t\) (\%)\\
\midrule
0.5 & 0.6034 & 0.2603 & 69.9 & 26.7\\
1.0 & 0.6835 & 0.2534 & 92.2 & 1.3\\
2.0 & 0.6770 & 0.2514 & 91.2 & 1.6\\
\bottomrule
\end{tabular}
\end{table}

\subsection{Control-basis dependence}

Replacing the Fourier basis with a modified-Chebyshev parameterization while holding
all other conditions identical produces a qualitative reversal: \(f_2\le6.2\%\)
at all three Reynolds numbers, while \(f_t>93\%\)
(Table~\ref{tab:basis_comparison}, Figure~\ref{fig:parametrization_dependence}).
This qualitative reversal upon changing only the parameterization basis demonstrates
that modal concentration is conditional on the chosen control representation,
not a universal property of the controlled cavity flow. The Chebyshev controls
also produce a higher realized energy (\(E_r\approx0.285\)--\(0.291\) vs.\
\(0.254\)--\(0.269\)), suggesting less efficient satisfaction of the energy constraint.
Basis construction details are provided in Supplementary Sec.~S7.

\begin{table*}[t]
\centering
\caption{Fourier versus modified-Chebyshev parameterization at \(E^*=0.25\).
All other conditions are identical.}
\label{tab:basis_comparison}
\small
\begin{tabular}{ccccccc}
\toprule
\(Re\) & \multicolumn{3}{c}{Fourier basis} & \multicolumn{3}{c}{Modified-Chebyshev basis}\\
\cmidrule(lr){2-4}\cmidrule(lr){5-7}
& \(E_r\) & \(f_2\) (\%) & \(f_t\) (\%) & \(E_r\) & \(f_2\) (\%) & \(f_t\) (\%)\\
\midrule
100  & 0.2581 & 91.6 & 0.7  & 0.2890 & 0.3 & 99.3\\
500  & 0.2534 & 92.2 & 1.3  & 0.2850 & 6.2 & 93.3\\
1000 & 0.2687 & 64.0 & 32.2 & 0.2908 & 0.5 & 98.8\\
\bottomrule
\end{tabular}
\end{table*}

\begin{figure*}[t]
\centering
\includegraphics[width=0.88\textwidth]{fig10_parametrization_dependence.pdf}
\caption{Control-basis dependence (Fourier vs.\ modified-Chebyshev) across three
Reynolds numbers at \(E^*=0.25\). Left: second-harmonic fraction \(f_2\). Right:
time-dependent fraction \(f_t\). The Fourier basis consistently selects predominantly
stationary controls; the Chebyshev basis selects time-dependent controls, demonstrating
that modal concentration is representation-conditional.}
\label{fig:parametrization_dependence}
\end{figure*}

\subsection{Multi-seed audit and label revalidation}

A dedicated audit campaign re-ran 99 PINN realizations with the final code across
18 configurations (5 seeds per configuration for the standard campaign cells, 8
for the over-sampled and large-capacity cells), covering the full range of branch
structures encountered in the main study. Every audited realization reproduces
the frozen branch label of the corresponding original run: configurations labeled
\(B_1\) such as cap\_large and cap\_small re-converge to \(B_1\) in 8/8
realizations; time-dominated configurations such as the modified-Chebyshev cells
(bas\_cheb) recover \(B_2\) in 5/5 realizations at every tested Reynolds number;
and the mixed cells reproduce their expected fractions (e.g.\ mc\_8x9:
\(B_1\)/\(B_2\)/\(B_3=2/2/1\); bas\_fourier\_re100: 1/3/1), with
ar\_rect\_half re-converging to the temporally dominated class (1/4/0). Independently, re-classifying the entire 185-realization reference campaign
(\(\lambda_{\mathrm{var}}=1\)) with
the four-class threshold rule (\(B_{\mathrm{first}}\): \(f_1>0.5\); \(B_1\): \(f_2>2f_t\) and \(f_2>0.5\); \(B_2\):
\(f_t>2f_2\) and \(f_t>0.5\); otherwise \(B_3\)) reproduces the frozen labels
with \textbf{zero disagreements (0/185)}, confirming the internal consistency of
the branch taxonomy and of the multi-seed audit. The per-configuration audit
summary is reported in Supplementary Sec.~S10.

%--------------------------------------------------------------
\section{INDEPENDENT MRT-LBM VERIFICATION AND NUMERICAL UNCERTAINTY}
\label{sec:lbm}
%--------------------------------------------------------------

\subsection{MRT-LBM solver and benchmark verification}

An independent two-dimensional nine-velocity (D2Q9) multiple-relaxation-time
lattice-Boltzmann (MRT-LBM) solver was used to evaluate selected PINN-identified
controls in a completely independent computational environment. The solver employs the D2Q9 lattice with the
MRT collision rule
\begin{equation}
\mathbf{m}^*=\mathbf{m}-\mathbf{S}(\mathbf{m}-\mathbf{m}^{\mathrm{eq}}),
\quad\mathbf{f}=\mathbf{M}^{-1}\mathbf{m}^*,
\end{equation}
where \(\mathbf{m}=\mathbf{M}\mathbf{f}\) are moment-space populations,
\(\mathbf{S}=\mathrm{diag}(s_0,\ldots,s_8)\) is the diagonal relaxation matrix,
and \(\mathbf{m}^{\mathrm{eq}}\) are the equilibrium moments. Lid boundary conditions
are imposed using the Zou--He velocity bounce-back scheme~\cite{pan2006high,chen1998lattice}.
The LBM implementation is entirely independent of the PINN framework: it uses a
separate Python/Numba solver (parallelized over the lattice with Numba
\texttt{prange}), a different spatial discretization (Cartesian lattice
vs.\ scattered collocation), and a fundamentally different numerical method.
Complete solver specifications are provided in Supplementary Sec.~S8.

The solver was validated against the benchmark cavity solutions of Ghia et al.\
\cite{ghia1982high}. Centerline velocity errors are reported in
Table~\ref{tab:ghia_verification}. At \(Re=100\) the \(L_2\) error in
\(u\) is \(0.00506\) and the \(L_\infty\) error is \(0.01283\). At \(Re=1000\)
the errors increase moderately to \(L_2(u)=0.01411\) and \(L_\infty(u)=0.03003\),
consistent with finer velocity gradient structure and the \(256\times256\) lattice
used for benchmark comparison. These errors are within typical MRT-LBM accuracy
at this resolution and confirm solver reliability for the controlled-flow study.

\begin{table}[t]
\centering
\caption{MRT-LBM centerline velocity errors versus Ghia et al.\
\protect\cite{ghia1982high} at reference grid \(N=256\).}
\label{tab:ghia_verification}
\small
\begin{tabular}{ccccc}
\toprule
\(Re\) & \(L_2(u)\) & \(L_\infty(u)\) & \(L_2(v)\) & \(L_\infty(v)\)\\
\midrule
100  & 0.00506 & 0.01283 & 0.00675 & 0.01136\\
1000 & 0.01411 & 0.03003 & 0.01321 & 0.02125\\
\bottomrule
\end{tabular}
\end{table}

\subsection{Controlled-flow dissipation comparison}

Four controlled configurations were evaluated at \(Re=500\) and \(Re=1000\):
(i) uniform lid (\(U_{\mathrm{lid}}=1\), the baseline); (ii) first-harmonic
\(\sin(\pi x)\); (iii) stationary second-harmonic \(A_2\sin(2\pi x)\) representing
the \(B_1\) branch; and (iv) the mean PINN control evaluated at \(t=0\),
representing the ensemble-averaged physics-informed solution. Integrated viscous
dissipation \(\varepsilon\) and percentage reduction versus uniform lid are
reported in Table~\ref{tab:lbm_dissipation}.

\begin{table*}[t]
\centering
\caption{Independent MRT-LBM integrated dissipation \(\varepsilon\) and percentage
reduction versus the uniform-lid reference for four controlled configurations.}
\label{tab:lbm_dissipation}
\small
\begin{tabular}{lcccc}
\toprule
& \multicolumn{2}{c}{\(Re=500\)} & \multicolumn{2}{c}{\(Re=1000\)}\\
\cmidrule(lr){2-3}\cmidrule(lr){4-5}
Configuration & \(\varepsilon\) & Reduction (\%) & \(\varepsilon\) & Reduction (\%)\\
\midrule
Uniform lid         & 0.0442824  & ---  & 0.0266676  & ---\\
\(\sin(\pi x)\)     & 0.0156278  & 64.7 & 0.0101974  & 61.8\\
\(A_2\sin(2\pi x)\) & 0.0104070  & 76.5 & 0.00472797 & 82.3\\
PINN (mean)         & 0.0118915  & 73.1 & 0.00524634 & 80.3\\
\bottomrule
\end{tabular}
\end{table*}

Both the stationary second-harmonic control \(A_2\sin(2\pi x)\) and the PINN mean
control achieve reductions of approximately \(73\)--\(83\%\) in integrated viscous
dissipation relative to the uniform-lid baseline, confirming that these controls
produce genuine physical improvements reproducible in an independent solver.
The close dissipation values of \(A_2\sin(2\pi x)\) and the PINN mean control
are particularly significant: the difference at \(Re=500\) is approximately
\(3.3\%\) relative to the uniform baseline, indicating that the low-dimensional
second-harmonic structure captures the principal low-dissipation character of
the learned PINN control. This provides independent physical evidence that the
modal concentration observed in the PINN optimization is connected to a genuine
dissipation mechanism in the controlled cavity.

\subsection{Grid-convergence uncertainty analysis}

A systematic three-grid refinement study was conducted for \(A_2\sin(2\pi x)\)
at \(Re=500\) using grid resolutions \(N\in\{128,256,512\}\). The resulting
dissipation values are \(\varepsilon_{128}=0.00935\), \(\varepsilon_{256}=0.010407\),
and \(\varepsilon_{512}=0.010966\). The observed order of convergence is
\(p_{\mathrm{obs}}=0.92\), and the Richardson-extrapolated value is
\(\varepsilon_\infty=0.01159\). The grid-convergence indices (GCI) are
\(\mathrm{GCI}_{256}=14.2\%\) and \(\mathrm{GCI}_{512}=7.1\%\). The uncertainty
is non-negligible for absolute values but small compared to the \(76.5\%\) reduction
relative to the uniform lid (76.5\% at \(N=256\)) and remains at 73.8\% when
evaluated against the Richardson-extrapolated value \(\varepsilon_\infty=0.01159\),
both far outside the GCI band: the dissipation conclusion is robust with respect
to the LBM discretization. Figure~\ref{fig:gci_updated} shows the convergence
behavior, and the full GCI estimate is detailed in Supplementary Sec.~S8.

The observed order \(p_{\mathrm{obs}}=0.92\) is the \emph{empirical} order of
convergence of the integrated observable \(\varepsilon\) for this specific
boundary condition, not the formal order of the MRT-LBM scheme itself (which
is nominally second-order in space). Several features of the controlled cavity
depress the apparent order and must be acknowledged when interpreting the GCI
values. First, the moving lid enters the calculation through a velocity boundary
condition imposed at every iteration: the discontinuous corner conditions generate
periodic grid-scale disturbances whose contribution to the strain-rate field---and
therefore to the integrated dissipation, a derivative-squared functional---does not
vanish cleanly under mesh refinement. Second, \(\varepsilon\) is dominated by the
near-wall velocity gradients; the first-order one-sided finite differences used
for \(\tau_w\) together with the under-resolved boundary layer at moderate \(N\)
inflate the absolute values, which is precisely why the extrapolated
\(\varepsilon_\infty=0.01159\) lies \emph{above} the \(N=128\) value \((0.00935)\)
while the sequence is non-monotone in \(N\). Third, the GCI is an indicator of
resolution adequacy, not a proof of second-order convergence:
\(p_{\mathrm{obs}}\approx0.92\) reflects the contribution of boundary and corner
artifacts to the integrated observable; consequently \(\mathrm{GCI}_{256}=14.2\%\)
means that the absolute dissipation at \(N=256\) carries roughly a 14\%
grid-uncertainty band for this observable.

The conclusions of this study are insensitive to this band: (i) relative
comparisons are made at fixed \(N\), so the uncertainty largely cancels in the
ratios (e.g.\ \(\varepsilon_{k=1}/\varepsilon_{\mathrm{unif}}\) at fixed
\(N=512\)); (ii) the dissipation reduction of the second-harmonic control relative to the
uniform lid (a non-iso-energetic comparison) is \(76.5\%\) at \(N=256\) and
\(73.8\%\) using the extrapolated \(\varepsilon_\infty\), both far outside the
14\% band, so the qualitative conclusion of reduced dissipation relative to the
uniform baseline is unchanged, and the fixed-energy ranking
\(k=1<\,\mathrm{uniform}<\,k=2\) (Sec.~\ref{sec:focused_comparison}) is likewise
unaffected since it rests on ratios at a common grid; (iii) the same hierarchy
is reproduced at the coarser grid in the energy-matched comparison, providing a
resolution cross-check in addition to the Richardson estimate.

\begin{figure*}[t]
\centering
\includegraphics[width=0.88\textwidth]{fig12_gci_updated.pdf}
\caption{Three-grid convergence analysis for integrated LBM dissipation of the
\(A_2\sin(2\pi x)\) control at \(Re=500\). Left: dissipation \(\varepsilon\)
versus grid resolution \(N\), with Richardson extrapolation \(\varepsilon_\infty=0.01159\).
Right: absolute deviation from extrapolated value as a percentage.
\(\mathrm{GCI}_{256}=14.2\%\), \(\mathrm{GCI}_{512}=7.1\%\).}
\label{fig:gci_updated}
\end{figure*}

\subsection{Quasi-stationarity of controlled flows}

The temporal stability of controlled LBM flows is characterized by the kinetic-energy
fluctuation ratio \(R_K=K_{\mathrm{fluct}}/K_{\mathrm{mean}}\)
(Table~\ref{tab:quasistationarity}, Figure~\ref{fig:quasistationarity}).
Fourier-basis controls yield \(R_K\le0.65\%\) across all tested Reynolds numbers,
well below the 1\% quasi-stationarity threshold, confirming that the predominantly
stationary second-harmonic Fourier control establishes a quasi-steady flow state.
By contrast, modified-Chebyshev controls produce profoundly unsteady responses
with \(R_K\ge44\%\) at all Reynolds numbers and \(R_K=847\%\) at \(Re=100\).
The two representations therefore generate qualitatively different physical flow
states: Fourier \(B_1\) controls produce steady cavity flows with substantially
lower integrated dissipation than the uniform-lid baseline, while Chebyshev
controls produce highly oscillatory, energetically costly flows. (At a fixed
actuation energy the second-harmonic control is, however, more dissipative than
the first harmonic and the uniform lid; Sec.~\ref{sec:focused_comparison}.)
The quasi-stationarity diagnostic provides independent physical confirmation that
the Fourier \(B_1\) branch corresponds to a physically realizable steady-state
control target.

Statistical stationarity of the temporally dominated \(B_2\) branch was additionally
verified at \(Re=500\), \(E^*=0.25\) over two long LBM windows
(\(T=128\) actuation periods and \(2T=256\) periods): the kinetic-energy
fluctuation ratio is \(R_K=14.9594\) in both windows, with a relative difference
\(|\Delta R_K|/R_K(2T)=1.6\times10^{-10}\), far below the 15\% persistence
threshold. The \(B_2\) flow is therefore statistically stationary in the
mean-dissipation sense even though it is strongly time-dependent at the actuation
frequency (\(R_K\simeq15\), confirming a genuinely unsteady regime clearly
distinct from the quasi-steady Fourier \(B_1\) state), while ruling out a
statistically evolving state. Window-level statistics are reported in
Supplementary Sec.~S11.

\begin{table}[t]
\centering
\caption{LBM kinetic-energy fluctuation ratio \(R_K\) (\%) for Fourier and
modified-Chebyshev controls. Values below 1\% indicate quasi-stationary flow.}
\label{tab:quasistationarity}
\small
\begin{tabular}{ccc}
\toprule
\(Re\) & Fourier \(R_K\) (\%) & Modified-Chebyshev \(R_K\) (\%)\\
\midrule
100  & 0.15 & 847\\
500  & 0.04 & 44\\
1000 & 0.65 & 125\\
\bottomrule
\end{tabular}
\end{table}

\begin{figure*}[t]
\centering
\includegraphics[width=0.88\textwidth]{fig13_quasistationarity.pdf}
\caption{Quasi-stationarity diagnostic for Fourier and modified-Chebyshev controls.
Left: time history of domain-integrated kinetic energy at \(Re=500\). Right:
fluctuation ratio \(R_K\) at all tested Reynolds numbers. Fourier controls remain
below the 1\% threshold; Chebyshev controls produce highly unsteady flows.}
\label{fig:quasistationarity}
\end{figure*}

%--------------------------------------------------------------
\subsection{Focused fixed-energy comparison of competing control classes}
\label{sec:focused_comparison}

To isolate the physical consequences of the competing control branches, a focused
case study is performed at fixed Reynolds number (\(Re=500\)) and fixed actuation energy
budget (\(E^*=0.25\)). We compare the independent LBM dissipation of four specific
actuation profiles matched precisely to this budget: (i) the uniform lid reference
(\(U_{\mathrm{lid}} = \sqrt{E^*} = 0.5\)), (ii) a pure stationary first-harmonic
control (\(A_1\sin(\pi x)\) with \(A_1 = \sqrt{2E^*} = 0.707\)), (iii) a pure stationary
second-harmonic control (\(A_2\sin(2\pi x)\) with \(A_2 = 0.707\)) representative
of the \(B_1\) branch, and (iv) representative temporally dominated (\(B_2\)) and
mixed (\(B_3\)) realizations drawn from the PINN campaign.

When enforcing a strict iso-energetic comparison at \(E^*=0.25\) on the two tested
grid resolutions (\(N=256\) and \(N=512\)), the second-harmonic \(B_1\) profile is
found to be \emph{more} dissipative than the uniform-lid baseline at both
resolutions, with a dissipation ratio
\(\varepsilon_{B_1}/\varepsilon_{\mathrm{uniform}}\) in the range
\(\left[1.072\;(N=512),\;1.188\;(N=256)\right]\).
In contrast, the first-harmonic profile (\(k=1\)) reduces dissipation below the
uniform baseline at the same energy budget at both resolutions, with
\(\varepsilon_{k=1}/\varepsilon_{\mathrm{uniform}}\) in the range
\(\left[0.646\;(N=512),\;0.726\;(N=256)\right]\).
The energy-matched ranking $k=1$\,$<$\,uniform\,$<$\,$k=2$ ($B_1$) is therefore
unchanged across the tested grids. The pure first-harmonic profile (ii) is
precisely the \(B_{\mathrm{first}}\) branch: the branch that is *not* selected at
the reference weight \(\lambda_{\mathrm{var}}=1\) is the one that is energetically
preferable to both the uniform lid and \(B_1\). This reinforces that branch
accessibility, not physical optimality, is what the reference loss formulation
controls.
However, the PINN optimizer systematically selects the second-harmonic \(B_1\)
structure over the first-harmonic structure under the variance-regularization
loss \(\mathcal{L}_{\mathrm{var}}\). Under the present sine-basis parameterization
this regularizer favors sign-changing controls with larger spatial variance rather
than penalizing wavenumber directly:
\(\operatorname{Var}[\sin(\pi x)]=1/2-4/\pi^2\approx0.0947\) for the positive
single-lobe profile versus \(\operatorname{Var}[\sin(2\pi x)]=1/2\) for the
balanced two-lobe profile. The term therefore biases branch selection toward the
second harmonic over the positive first-harmonic profile at matched actuation
energy.

The temporal \(B_2\) and mixed \(B_3\) controls produce strongly unsteady flows
(for the \(B_2\) reference realization the kinetic-energy fluctuation ratio is
\(R_K\simeq14.96\), reproduced identically over both \(T\) and \(2T\) statistical
windows), breaking the quasi-stationarity required for direct mean-dissipation
comparisons. These results emphasize a central conclusion of this study: the
frequent selection of the \(B_1\) branch by the PINN optimizer does not imply
that \(B_1\) is the global physical optimum for dissipation at a fixed energy budget.
Instead, it demonstrates that branch accessibility is jointly determined by the
physical constraints and the specific regularizations of the loss landscape
(branch accessibility \(\neq\) global optimality). Full methodological details and
energy-matched sweep data are provided in Supplementary Sec.~S11.

\section{DISCUSSION: CONDITIONAL MODAL BRANCH SELECTION}
\label{sec:discussion}
%--------------------------------------------------------------

\subsection{What is actually selected by the physics-informed optimization?}

The experiments presented above establish that the output of the physics-informed
optimization is not a deterministic map \((Re,E^*)\to\mathcal{B}\) but rather
an empirical frequency distribution over accessible branch classes that depends
on the full set of variational parameters:
\(\mathcal{B}(Re,E^*,\lambda_{\mathrm{var}},\text{seed},\text{capacity},\text{basis},AR)\).
This finding has a direct methodological implication: a single converged solution
can only be said to be a solution that the adopted optimization procedure can reach;
it cannot, without additional multi-start evidence, be said to be the unique or
globally optimal solution of the underlying control problem.

The co-existence of \(B_1\) and \(B_2\) realizations with statistically
indistinguishable total losses (spread \(<2\%\)) is the most direct evidence for
the competing-basin interpretation. If one branch were substantially more optimal,
it would achieve measurably lower losses, which is not observed. The near-equal
global accessibility of \(B_1\) (42.2\%) and \(B_2\) (43.8\%) further confirms
that neither branch dominates globally. The concept of a ``dominant'' mode must
therefore be replaced by ``frequently accessible under the given formulation.''

The \(B_3\) mixed branch, accounting for 14.1\% of realizations in the reference
campaign, is selectively favored by high actuation energy and elevated Reynolds
number; its less frequent appearance under standard loss weights suggests it lies
in shallower basins. The first-harmonic branch \(B_{\mathrm{first}}\), absent at the
reference weight \(\lambda_{\mathrm{var}}=1\), becomes the exclusive outcome once
the spatial-variance penalty is removed (\(\lambda_{\mathrm{var}}=0\), \(H_N=0\)):
without this penalty the optimizer selects the lowest-order stationary mode, the
single-lobe profile \(\sin(\pi x)\), precisely the control that the variance
regularizer disfavors (spatial variance \(0.0947\) versus \(1/2\)) in favour of the
two-lobe second-harmonic structure. The two
branches \(B_{\mathrm{first}}\) and \(B_1\) therefore bracket the range of spatially
stationary controls that the variational formulation can select, and the choice
between them is controlled by a single loss weight rather than by the flow physics.

\subsection{Physical origin of second-harmonic recurrence}

Despite the caveats about accessibility and global optimality, the second-harmonic
branch \(B_1\) has a clear physical rationale. The antisymmetric profile
\(A_2\sin(2\pi x)\) creates a two-lobe actuation pattern at the upper lid,
generating a pair of counter-rotating recirculation cells rather than the single
primary vortex of the uniform-lid cavity. This produces a more symmetric and
geometrically balanced velocity gradient distribution, reducing the regions of
intense shear that dominate viscous dissipation. The resulting \(70\)--\(82\%\)
dissipation reduction relative to the uniform-lid baseline, confirmed by the
independent LBM calculations, provides direct physical validation of the modal
structure under a non-iso-energetic comparison. At a fixed actuation energy the
second-harmonic control is in fact more dissipative than both the first harmonic
and the uniform lid (Sec.~\ref{sec:focused_comparison}); its frequent selection
therefore reflects accessibility under the optimization formulation rather than
energetic optimality.

The quasi-stationarity of the Fourier \(B_1\) controlled flow (\(R_K\le0.65\%\))
is a further physical confirmation of \emph{realizability}: the actuation profile
selected by the optimizer generates a steady cavity state, showing that the branch
is attainable as a physically realizable steady control. Realizability is distinct
from energetic optimality: at a fixed actuation energy the iso-energetic comparison
of Sec.~\ref{sec:focused_comparison} shows \(B_1\) to be more dissipative than both
the first harmonic and the uniform lid. In contrast, the Chebyshev controls produce
highly unsteady flows (\(R_K\ge44\%\)), representing a class of control response in
which added kinetic energy of temporal oscillations prevents convergence to a
quasi-stationary steady state. This qualitative physical difference
provides an independent explanation for why the two representations lead to
different branch outcomes.

The persistence of the second-harmonic structure across capacities and truncations
is consistent with it being a genuine physical steady-state mode: it yields
quasi-stationary flows with dissipation far below the uniform-lid baseline
(Table~\ref{tab:lbm_dissipation}). The fixed-energy ranking of
Sec.~\ref{sec:focused_comparison} nevertheless shows that it is not the
least-dissipative stationary profile at a common actuation budget. However,
the erosion of \(f_2\) at high energy (\(E^*=2\)) and at compressed geometry
(\(AR=0.5\)) shows that the second-harmonic basin is not infinitely deep: when
the energy budget or geometry changes sufficiently, other modal strategies become
competitive. The global balance between \(B_1\) and \(B_2\) accessibility in the
broad campaign reflects this fundamental competitiveness between stationary and
time-dependent dissipation strategies.

\subsection{Optimization bias versus fluid mechanics}

The loss-ablation results demonstrate that the dominant driver of \(B_1\) selection
is \(\mathcal{L}_{\mathrm{var}}\), not \(\mathcal{L}_{\mathrm{diss}}\). Removing
\(\mathcal{L}_{\mathrm{diss}}\) leaves \(f_2\) essentially unchanged (\(92.3\%\)
vs.\ baseline \(92.2\%\)), while removing \(\mathcal{L}_{\mathrm{var}}\) causes a
complete modal-class shift to the first-harmonic branch \(B_{\mathrm{first}}\)
(\(f_1=90.4\%\), \(f_2=0.2\%\)). This means that in the reference loss
formulation, the spatial-variance penalty acts as the primary optimization landscape
shaper, directing the search into the basin of attraction of the spatially structured
\(B_1\) control. The physical dissipation objective, while present and contributing
to overall convergence quality, plays a secondary role in determining which basin
is entered.

This does not invalidate the physical relevance of the identified controls: the
LBM verification confirms that the controls reduce integrated viscous dissipation
relative to the uniform-lid baseline, and the quasi-stationarity analysis confirms
they are physically realizable. What the ablation result reveals is a distinction
between \emph{which loss term selects the branch} and \emph{which physical property
justifies the selected branch}. The variance regularization acts as an implicit
selector of controls with larger spatial variance (in the sine basis, the
sign-changing two-lobe profile over the positive single-lobe one); the dissipation
objective then evaluates those controls. They reduce dissipation relative to the
uniform lid, but under a strictly iso-energetic comparison the first harmonic is
the least dissipative of the tested stationary profiles
(Sec.~\ref{sec:focused_comparison}). Practitioners
designing physics-informed control problems should be aware that the loss weight
configuration may be as consequential as the physical objective in determining which
branch is selected.

\subsection{Scope, limitations, and directions for extension}

The present study is bounded to two-dimensional incompressible flow, quasi-stationary
control objectives, and a Fourier actuation basis with fixed spatial support on the
upper lid. Several directions of extension present themselves. First, extension to
three-dimensional cavity flow would introduce additional modal complexity, including
spanwise modes and potential Taylor-G\"{o}rtler vortex structures. Second, an
explicitly unsteady objective might access different branches and provide a more
complete picture of the control landscape. Third, higher-fidelity Navier--Stokes
solvers at higher resolution would reduce the GCI uncertainty and provide sharper
bounds on absolute dissipation values.

The representation-dependence result also raises the broader question of how other
basis choices (Zernike polynomials, wavelets, data-driven POD modes) would partition
the accessible branch space. This question cannot be answered by physical arguments
alone but requires systematic multi-representation campaigns of the type conducted
here. The symbolic-regression analysis in the following section provides a compact
algebraic summary that may serve as a starting point for analytical investigation
of the control mechanism at larger Reynolds numbers.

%--------------------------------------------------------------
\section{COMPACT REPRESENTATION BY SYMBOLIC REGRESSION}
\label{sec:sr}
%--------------------------------------------------------------

\subsection{Symbolic regression of the optimal-control profile}

The \(B_1\) Fourier controls share a common spatial structure that is well
approximated by a compact analytical expression. We apply the PySR symbolic
regression framework~\cite{cranmer2023pysr} to a dataset of \(B_1\) control
coefficients, seeking an algebraic expression for the mean lid velocity profile
\(\bar U_{\mathrm{lid}}(x;Re)\). PySR searches over binary operators
\(\{+,-,\times,/\}\) and unary operators \(\{\sin,\cos,\exp,\mathrm{square}\}\)
using an evolutionary algorithm across 20 independent populations with 100
optimization iterations each. The best complexity-accuracy trade-off on the
Pareto frontier yields
\begin{equation}
U_{\mathrm{SR}}(x;Re)\approx
a\left[b+\cos(\omega_{\mathrm{SR}}x-\phi)\right]^2
(c_0+c_1 x+c_2 t+c_3\,Re)+c_4,
\quad\omega_{\mathrm{SR}}\approx 3.42.
\label{eq:sr}
\end{equation}
The spatial frequency \(\omega_{\mathrm{SR}}\approx3.42\) is close to \(\pi\approx3.14\),
consistent with the squared cosine structure generating an effective second-harmonic
pattern. The squared cosine structure \([b+\cos(\omega_{\mathrm{SR}}x-\phi)]^2\)
provides a smooth, non-negative envelope capturing the two-lobe velocity pattern
of the \(B_1\) branch. The linear dependence on \(Re\) reflects the moderate drift
of optimal amplitude with Reynolds number. The full PySR configuration and dataset
specification are provided in Supplementary Sec.~S9.

\subsection{Leave-one-Reynolds-condition-out validation}

A leave-one-Reynolds-condition-out (LOCO) cross-validation is performed across 9
held-out Reynolds values \(Re\in\{100,300,400,500,600,700,800,900,1000\}\).
For each held-out condition, PySR is re-fit on the remaining 8 conditions, and
the resulting expression is evaluated on the held-out data. Figure~\ref{fig:symbolic_loco}
shows the LOCO \(R^2\) scores. The scores range from \(R^2=0.9326\) (worst case,
\(Re=900\)) to \(R^2=0.9908\) (best case, \(Re=500\)), with all 9 values exceeding
\(0.93\). This consistent fidelity demonstrates that the symbolic expression
captures a genuine cross-condition structure of the \(B_1\) control family, not
a case-specific interpolation artifact. Condition-wise LOCO scores and the full
PySR configuration are provided in Supplementary Sec.~S9.

\begin{figure}[t]
\centering
\includegraphics[width=0.95\columnwidth]{fig14_symbolic_loco.pdf}
\caption{Leave-one-Reynolds-condition-out (LOCO) \(R^2\) cross-validation scores
for the symbolic regression expression Eq.~\eqref{eq:sr}. All scores exceed
\(0.93\), confirming robust cross-condition generalization of the identified
\(B_1\) control family.}
\label{fig:symbolic_loco}
\end{figure}

\subsection{Scope and interpretation of the symbolic result}

The symbolic expression Eq.~\eqref{eq:sr} provides three useful contributions:
it demonstrates that the \(B_1\) control family can be compactly described by a
smooth algebraic function with interpretable parameters; the cross-condition
fidelity confirms that \(B_1\) is a coherent family of related solutions varying
smoothly with \(Re\); and the expression provides a compact, training-free
approximation to the optimal Fourier profile suitable as an initial guess in future
control studies. However, the symbolic expression is a compact data description,
not an independent proof of optimality. High LOCO \(R^2\) values establish that the
branch structure is internally consistent, not that the branch is globally optimal.

%--------------------------------------------------------------
\section{CONCLUSIONS}
\label{sec:conclusions}
%--------------------------------------------------------------

This study established a quantitative framework for investigating conditional modal
branch selection in physics-informed active flow control of the lid-driven cavity.
The principal findings are as follows.

\begin{enumerate}
\item \textbf{Four empirical branch classes are reproducibly identified.}
A broad \((Re,E^*)\) campaign of 185 independent PINN optimizations at the
reference weight \(\lambda_{\mathrm{var}}=1\) reveals three structurally distinct and
well-separated populated branch classes: \(B_1\) (stationary
second-harmonic dominated, 42.2\%), \(B_2\) (temporally dominated, 43.8\%), and
\(B_3\) (mixed, 14.1\%). A fourth class, the stationary first-harmonic branch
\(B_{\mathrm{first}}\), is not populated at this weight but emerges as the exclusive
outcome when the variance regularizer is removed (Sec.~\ref{sec:lambda_var}). The
near balance of \(B_1\) and \(B_2\) demonstrates
that neither branch uniquely dominates the optimization landscape; both are
frequently accessible.

\item \textbf{Branch accessibility depends on the optimization formulation.}
The variance-regularization weight \(\lambda_{\mathrm{var}}\) is the dominant
determinant of branch selection. At \(\lambda_{\mathrm{var}}=0\), all conditions
yield exclusively \(B_{\mathrm{first}}\) outcomes. Suppression of the first-harmonic
branch by a positive \(\lambda_{\mathrm{var}}\) is operating-point dependent: complete
already at \(\lambda_{\mathrm{var}}=0.25\) for \((Re,E^*)=(500,0.25)\) and
\((1000,1)\), but reached only at \(\lambda_{\mathrm{var}}=0.5\) for \((700,2)\).
Increasing \(\lambda_{\mathrm{var}}\) further non-monotonically shifts accessibility
between \(B_1\) and \(B_2\) in a way that depends on the operating point. The branch
structure is jointly a property of
the fluid problem and the variational formulation, not of the fluid problem alone.

\item \textbf{Random initialization produces competing branches of comparable
loss quality.}
Ten independent seed experiments at \(Re=500\), \(E^*=0.25\) produce 4 \(B_1\),
4 \(B_2\), and 2 \(B_3\) realizations with total training losses clustered within
\([1.371,1.394]\), a relative spread of only 1.7\%. The competing branches
correspond to basins of comparable depth, not to a unique global minimum.

\item \textbf{The stationary second-harmonic branch is robust to network capacity
and control-space dimension.}
The \(B_1\) branch persists across all tested network capacities (\(4\times32\)
to \(8\times128\)) and Fourier basis truncations (12 to 72 nominal modes),
with \(f_2\in[87.3\%,92.3\%]\). This robustness argues against the hypothesis
that the structure is a numerical artifact of insufficient expressivity.

\item \textbf{Modal concentration is representation-conditional.}
Replacing the Fourier basis with a modified-Chebyshev parameterization reverses
the modal outcome: Chebyshev controls are exclusively time-dependent (\(f_t>93\%\))
and produce highly unsteady cavity flows (\(R_K\ge44\%\)), while Fourier \(B_1\)
controls establish quasi-stationary flows (\(R_K\le0.65\%\)).

\item \textbf{Selected controls are independently verified and produce genuine
dissipation reductions.}
Transfer of the \(A_2\sin(2\pi x)\) and PINN mean controls to an independent
MRT-LBM solver confirms reductions of \(\sim76.5\%\) at \(Re=500\) and
\(\sim82.3\%\) at \(Re=1000\) in integrated viscous dissipation relative to
the uniform-lid baseline. The close dissipation values of the analytic
second-harmonic and the full PINN mean control provide independent evidence
that the low-dimensional structure captures the principal low-dissipation character
of the learned control.

\item \textbf{Symbolic regression provides a compact, generalizing description.}
PySR symbolic regression yields a compact expression with spatial frequency
\(\omega_{\mathrm{SR}}\approx3.42\), achieving leave-one-Reynolds-condition-out
\(R^2\in[0.9326,0.9908]\) across all nine held-out conditions.
The consistently high LOCO scores confirm that the \(B_1\) branch forms a coherent
family of related controls varying smoothly and predictably with Reynolds number.

\item \textbf{Branch accessibility does not imply global physical optimality.}
A focused iso-energetic LBM comparison at \(E^*=0.25\) on two grid resolutions
(\(N=256\), \(N=512\)) demonstrates that the frequently selected
second-harmonic control (\(B_1\)) is more dissipative than the uniform-lid
baseline at the same energy budget, with a dissipation ratio in the range
\([1.072,1.188]\), whereas the first harmonic is consistently less dissipative,
with a ratio in the range \([0.646,0.726]\). The energy-matched ranking
$k=1$\,$<$\,uniform\,$<$\,$k=2$ (\(B_1\)) is unchanged across the tested
resolutions. The PINN selects the \(B_1\) branch that is sub-optimal from the
standpoint of fixed-energy dissipation because the variance regularization
heavily penalizes low spatial frequencies. Synthesizing evidence from over 300
independent PINN realizations across all campaigns (185 broad, 90 variance, and
sensitivity sweeps), this confirms that optimization-landscape regularizations
drive branch selection rather than absolute physical optimality.
\end{enumerate}

The principal conclusion is summarized by the following hierarchy:
\begin{equation}
\boxed{\text{branch existence}\;\neq\;\text{branch accessibility}
\;\neq\;\text{branch reproducibility}\;\neq\;\text{global optimality}.}
\label{eq:hierarchy}
\end{equation}
The second-harmonic \(B_1\) branch satisfies the first three properties---it exists
as a physically realizable flow state, is accessible under appropriate variational
formulations, and is reproducible across computational platforms---but the evidence
presented here does not establish global optimality. The co-existence of competitive
\(B_2\) solutions at essentially the same loss level precludes such a claim.
The framework developed here, combining multi-start PINN campaigns, systematic
sensitivity analyses, and independent LBM verification, provides a general methodology
for the rigorous characterization of branch selection in physics-informed flow
control, applicable beyond the cavity geometry studied.

%--------------------------------------------------------------
\section*{Supplementary Material}
%--------------------------------------------------------------
The supplementary material is provided as a separate online document (file type
``Supplementary Online Material''). It contains the complete implementation and
data details supporting the manuscript: (S1) neural-network architecture and
sampling details; (S2) modal-energy quadrature conventions; (S3) realization-level
data of the 185-run reference campaign (Supplementary Table~S1); (S4) the
\(\lambda_{\mathrm{var}}\) sensitivity campaign with 95\% Wilson-score confidence
intervals; (S5) energy-sweep and seed-multiplicity datasets; (S6) network-capacity
and control-truncation experiments; (S7) spectral-basis and aspect-ratio
construction details; (S8) MRT-LBM solver specification, Ghia benchmark tables,
and the grid-convergence (GCI) calculation including the Richardson-extrapolated
estimate; (S9) PySR configuration and leave-one-Reynolds-condition-out (LOCO)
scores; (S10) the reproducibility double-run tables (seeds 4/7/8) and the
99-realization multi-seed audit summary; and (S11) the energy-matched sweep data,
the Exp~D trajectory tests (autocorrelation, control-cycle, and modal-drift
analyses), and the long-window quasi-stationarity statistics of the \(B_2\) branch.

%--------------------------------------------------------------
\section*{Acknowledgment}
%--------------------------------------------------------------
The authors thank the High-Performance Computing facility of the Ecole Militaire
Polytechnique for access to the computational resources used in this study.

%--------------------------------------------------------------
\section*{Funding}
%--------------------------------------------------------------
The authors received no financial support for this research.

\section*{Conflicts of Interest}
The authors declare no competing interests.

\section*{Data and Code Availability}
Code, datasets, and LBM verification workflows are publicly available at:
\url{https://github.com/Ahmedalgeria778/pinn-lid-driven-cavity-collapse}.

%--------------------------------------------------------------
\begin{thebibliography}{99}

\bibitem{gad2000flow}
M.\ Gad-el-Hak,
\textit{Flow Control: Passive, Active, and Reactive Flow Management}
(Cambridge University Press, 2000).

\bibitem{brunton2015closed}
S.\ L.\ Brunton and B.\ R.\ Noack,
``Closed-loop turbulence control: Progress and challenges,''
Appl.\ Mech.\ Rev.\ \textbf{67}, 050801 (2015).

\bibitem{rowley2017model}
C.\ W.\ Rowley and S.\ T.\ Dawson,
``Model reduction for flow analysis and control,''
Annu.\ Rev.\ Fluid Mech.\ \textbf{49}, 387--417 (2017).

\bibitem{taira2017modal}
K.\ Taira et al.,
``Modal analysis of fluid flows: An overview,''
AIAA J.\ \textbf{55}, 4013--4041 (2017).

\bibitem{bergmann2005optimal}
M.\ Bergmann, L.\ Cordier, and J.-P.\ Brancher,
``Optimal rotary control of the cylinder wake using proper orthogonal
decomposition reduced-order model,''
Phys.\ Fluids \textbf{17}, 097101 (2005).

\bibitem{bewley2001dns}
T.\ R.\ Bewley, P.\ Moin, and R.\ Temam,
``DNS-based predictive control of turbulence: an optimal benchmark for
feedback algorithms,''
J.\ Fluid Mech.\ \textbf{447}, 179--225 (2001).

\bibitem{rabault2019drl}
J.\ Rabault et al.,
``Artificial neural networks trained through deep reinforcement learning
discover control strategies for active flow control,''
J.\ Fluid Mech.\ \textbf{865}, 281--302 (2019).

\bibitem{garnier2021drl}
P.\ Garnier et al.,
``A review on deep reinforcement learning for fluid mechanics,''
Comput.\ Fluids \textbf{225}, 104973 (2021).

\bibitem{ghia1982high}
U.\ Ghia, K.\ N.\ Ghia, and C.\ T.\ Shin,
``High-Re solutions for incompressible flow using the Navier--Stokes
equations and a multigrid method,''
J.\ Comput.\ Phys.\ \textbf{48}, 387--411 (1982).

\bibitem{raissi2019pinn}
M.\ Raissi, P.\ Perdikaris, and G.\ E.\ Karniadakis,
``Physics-informed neural networks: A deep learning framework for solving
forward and inverse problems involving nonlinear partial differential
equations,''
J.\ Comput.\ Phys.\ \textbf{378}, 686--707 (2019).

\bibitem{cai2021pinnreview}
S.\ Cai et al.,
``Physics-informed neural networks (PINNs) for fluid mechanics: A review,''
Acta Mech.\ Sin.\ \textbf{37}, 1727--1738 (2021).

\bibitem{duraisamy2019turbulence}
K.\ Duraisamy, G.\ Iaccarino, and H.\ Xiao,
``Turbulence modeling in the age of data,''
Annu.\ Rev.\ Fluid Mech.\ \textbf{51}, 357--377 (2019).

\bibitem{brenner2019perspective}
M.\ P.\ Brenner, J.\ D.\ Eldredge, and J.\ B.\ Freund,
``Perspective on machine learning for advancing fluid mechanics,''
Phys.\ Rev.\ Fluids \textbf{4}, 100501 (2019).

\bibitem{lozano2020machine}
A.\ Lozano-Dur\'{a}n and G.\ Arranz,
``Information-theoretic formulation of dynamical systems: causality,
modeling, and control,''
Phys.\ Rev.\ Res.\ \textbf{4}, 023195 (2022).

\bibitem{schmid2010dmd}
P.\ J.\ Schmid,
``Dynamic mode decomposition of numerical and experimental data,''
J.\ Fluid Mech.\ \textbf{656}, 5--28 (2010).

\bibitem{noack2003hierarchy}
B.\ R.\ Noack et al.,
``A hierarchy of low-dimensional models for the transient and
post-transient cylinder wake,''
J.\ Fluid Mech.\ \textbf{497}, 335--363 (2003).

\bibitem{bagheri2009global}
S.\ Bagheri et al.,
``Global stability of a jet in crossflow,''
J.\ Fluid Mech.\ \textbf{624}, 33--44 (2009).

\bibitem{farano2016subcritical}
M.\ Farano et al.,
``Subcritical transition scenarios via linear and nonlinear localized
optimal perturbations in plane Poiseuille flow,''
J.\ Fluid Mech.\ \textbf{802}, R1 (2016).

\bibitem{beneddine2016conditions}
S.\ Beneddine et al.,
``Conditions for validity of mean flow stability analysis,''
J.\ Fluid Mech.\ \textbf{798}, 485--504 (2016).

\bibitem{sipp2010global}
D.\ Sipp and A.\ Lebedev,
``Global stability of base and mean flows: a general approach and its
applications to cylinder and open cavity flows,''
J.\ Fluid Mech.\ \textbf{593}, 333--358 (2007).

\bibitem{luhar2014opposition}
M.\ Luhar, A.\ S.\ Sharma, and B.\ J.\ McKeon,
``Opposition control within the resolvent analysis framework,''
J.\ Fluid Mech.\ \textbf{749}, 597--626 (2014).

\bibitem{pan2006high}
X.\ F.\ Pan et al.,
``Lattice Boltzmann simulation of particle-laden turbulent channel flow,''
Comput.\ Fluids \textbf{35}, 1432--1444 (2006).

\bibitem{chen1998lattice}
S.\ Chen and G.\ D.\ Doolen,
``Lattice Boltzmann method for fluid flows,''
Annu.\ Rev.\ Fluid Mech.\ \textbf{30}, 329--364 (1998).

\bibitem{cranmer2023pysr}
M.\ Cranmer,
``Interpretable machine learning for physics,''
arXiv:2305.01582 (2023).

\bibitem{lai2000lattice}
Y.-G.\ Lai and R.\ S.\ Myong,
``A generalized wall-function treatment for high-Reynolds-number flows
over complex configurations,''
AIAA Paper 2000-0143 (2000).

\end{thebibliography}

\end{document}
